% Co-governed and enforced under the Sovereign Integrity Protocol License (SIP License v1.1):
% https://github.com/tensorrent/ACS-Framework/blob/main/LICENSE
\documentclass[11pt,a4paper]{article}

\usepackage[utf8]{inputenc}
\usepackage[T1]{fontenc}
\usepackage{amsmath,amssymb,amsthm,bm}
\usepackage{booktabs}
\usepackage{microtype}
\usepackage{xcolor}
\usepackage{hyperref}
\usepackage{enumitem}
\usepackage{geometry}
\geometry{margin=0.85in}

\hypersetup{
  colorlinks=true,
  linkcolor=blue!70!black,
  citecolor=green!50!black,
  urlcolor=red!70!black
}

\newtheorem{proposition}{Proposition}
\newtheorem{corollary}{Corollary}
\newtheorem{remark}{Remark}

\title{\textbf{The Framing Transformer:\\
Where the Spin-$1/2$ Content of the M\"obius-Screw Actually Lives}\\[0.35em]
\large Companion computation to \textit{The M\"obius-Screw Electron}}
\author{Flag Condensate Collaboration\\
\small\textit{Sovereign-Stack ACS Research Program}}
\date{July 26, 2026}

\begin{document}
\maketitle

\begin{abstract}
The M\"obius-screw note models the electron as a framed unknot with
self-linking number $Sl=2$ under the torus framing, and identifies
$Sl \leftrightarrow g$ as a bookkeeping statement.  This note computes
the full chain of maps from the centerline to the $SU(2)$ lift of its
frame field, and evaluates every stage.  Three results.
\emph{(i)} The C\u{a}lug\u{a}reanu decomposition is confirmed
numerically: $Tw = -1.0338$, $Wr = -0.9662$, $Tw+Wr = -2.000000$, and the
Gauss linking integral of the centerline with its framing pushoff returns
$-2.000000$; the magnitude $|Sl| = pq = 2$ is as stated, the sign records
that the embedding as parameterised is left-handed.
\emph{(ii)} The frame loop is genuinely \emph{spinorial}: its class in
$\pi_1(SO(3)) \cong \mathbb{Z}/2$ is non-trivial, so the $SU(2)$ lift
does not close after one circuit.  This is proved in closed form, not
merely computed.
\emph{(iii)} The spin content is carried by a \emph{parity bit}, not by
the value $2$.  For torus curves the holonomy is
$\sigma = (-1)^{p+q} = (-1)^{Sl+1}$, so $Sl = 2$ and $Sl = 0$ lie in the
same class: an untwisted round circle is spinorial in exactly the same
sense.  The magnitude of $Sl$ carries no spin information, which
withdraws the support the $Sl=2 \leftrightarrow g=2$ identification was
taken to have.  What survives is sharper and smaller: the double cover in
this model comes from the meridian half-angle $\varphi/2$ in the
parameterisation, i.e.\ from the \emph{odd} meridian winding $q=1$.
\end{abstract}

\tableofcontents
\newpage

\section{Scope}

This note is a computation, not a new model.  It takes the centerline and
framing already defined in the M\"obius-screw note and asks a single
question: \emph{at which stage, if any, does spin-$1/2$ enter?}  Every
number below is produced by
\texttt{code/framed\_unknot/framing\_transformer.py} and archived in
\texttt{docs/framed\_unknot\_results.json}.

Nothing here bears on QED, on radiative corrections to $g-2$, or on the
capacitance estimate of $\alpha$.  The result is topological.

\section{The transformer chain}

The object is the closed curve, $\varphi \in [0,4\pi)$,
\begin{equation}
\label{eq:curve}
\bm{\gamma}(\varphi) =
\begin{pmatrix}
\bigl(R + a\cos(\varphi/2)\bigr)\cos\varphi \\
\bigl(R + a\cos(\varphi/2)\bigr)\sin\varphi \\
a\sin(\varphi/2)
\end{pmatrix},
\end{equation}
which is the $(p,q) = (2,1)$ curve on the torus of radii $(R,a)$:
longitude angle $\theta = \varphi$, meridian angle $\psi = \varphi/2$.
The torus surface normal
\begin{equation}
\label{eq:framing}
\bm{U}(\varphi) = \bigl(\cos(\varphi/2)\cos\varphi,\;
                        \cos(\varphi/2)\sin\varphi,\;
                        \sin(\varphi/2)\bigr)
\end{equation}
is the natural framing.  Because $\bm{\gamma}$ lies in the surface whose
normal is $\bm{U}$, orthogonality $\bm{T}\cdot\bm{U} = 0$ holds
identically (verified to $2.2\times10^{-16}$).

The chain from geometry to spin is
\[
\bm{\gamma}
\;\xrightarrow{\ \text{framing}\ }\;
\bm{U}
\;\xrightarrow{\ \text{CWF}\ }\;
Sl = Tw + Wr
\;\xrightarrow{\ \text{adapted frame}\ }\;
F : S^1 \to SO(3)
\;\xrightarrow{\ \text{lift}\ }\;
q : S^1 \to SU(2),
\]
with $F(\varphi) = [\bm{T},\bm{U},\bm{T}\times\bm{U}]$.  The last arrow is
the only stage at which a double cover can appear, and whether it
succeeds is decided by the single sign
\begin{equation}
\label{eq:sigma}
\sigma \;=\; \frac{q(4\pi)}{q(0)} \in \{+1,-1\},
\end{equation}
$\sigma = +1$ meaning the lift closes (frame returns after one circuit)
and $\sigma = -1$ meaning it does not (two circuits required).

\section{Stage 1: the C\u{a}lug\u{a}reanu decomposition}

Twist is evaluated as
$Tw = \tfrac{1}{2\pi}\oint (\bm{T}\times\bm{U})\cdot \bm{U}'\,d\varphi$,
writhe by the Gauss double integral, and $Sl$ independently as the linking
number $\mathrm{Lk}(\bm{\gamma},\bm{\gamma}+\epsilon\bm{U})$.

\begin{center}
\begin{tabular}{@{}lrl@{}}
\toprule
Quantity & Value & Method \\
\midrule
$Tw$                 & $-1.033761$ & twist integral \\
$Wr$                 & $-0.966239$ & Gauss double integral \\
$Tw + Wr$            & $-2.000000$ & C\u{a}lug\u{a}reanu sum \\
$\mathrm{Lk}(\bm{\gamma},\bm{\gamma}+0.10a\bm{U})$ & $-2.000000$ & Gauss linking integral \\
$\mathrm{Lk}(\bm{\gamma},\bm{\gamma}+0.05a\bm{U})$ & $-2.000019$ & same, halved offset \\
$pq$                 & $2$         & torus-framing formula \\
\bottomrule
\end{tabular}
\end{center}

The magnitude $|Sl| = 2$ claimed in the M\"obius-screw note is confirmed
by two independent routes.  The negative sign is not an error: it records
that the embedding \eqref{eq:curve}, with meridian angle advancing as
$+\varphi/2$, is left-handed under the standard Gauss convention.  Only
$|Sl|$ and its parity are used below.

\begin{remark}[What is and is not topological here]
$Tw$ and $Wr$ individually depend on $R$ and $a$ and are not invariants;
only their sum is.  The near-equal split
$(-1.03, -0.97)$ at $a/R = 0.3$ is a coincidence of that aspect ratio, not
a structural fact, and should not be read as ``one unit of twist plus one
unit of writhe.''
\end{remark}

\section{Stage 2: the $SU(2)$ lift, in closed form}

The framing \eqref{eq:framing} is exactly the sphere parameterisation at
longitude $\varphi$ and latitude $\varphi/2$, so
$\bm{U}(\varphi) = A(\varphi)\,\bm{e}_x$ with
\begin{equation}
A(\varphi) = R_z(\varphi)\, R_y(-\varphi/2).
\end{equation}
Lifting each factor to unit quaternions
($R_z(t)\mapsto \cos\tfrac{t}{2} + k\sin\tfrac{t}{2}$,
 $R_y(t)\mapsto \cos\tfrac{t}{2} + j\sin\tfrac{t}{2}$) gives the
continuous lift with $q(0)=1$,
\begin{equation}
\label{eq:lift}
q(\varphi) =
\Bigl[\cos\tfrac{\varphi}{2} + k\,\sin\tfrac{\varphi}{2}\Bigr]
\Bigl[\cos\tfrac{\varphi}{4} - j\,\sin\tfrac{\varphi}{4}\Bigr].
\end{equation}

\begin{proposition}[Spinorial holonomy]
\label{prop:spinorial}
For the curve \eqref{eq:curve} with framing \eqref{eq:framing},
\[
q(4\pi) = \bigl[\cos 2\pi + k \sin 2\pi\bigr]\,
          \bigl[\cos \pi - j \sin \pi\bigr]
        = (1)(-1) = -1 ,
\]
so $\sigma = -1$: the frame loop represents the non-trivial class of
$\pi_1(SO(3)) \cong \mathbb{Z}/2$ and its $SU(2)$ lift does not close
after one circuit of the curve.
\end{proposition}

\begin{proof}[Where the sign comes from]
The meridian factor.  Over the closed curve the meridian angle advances
by $2\pi$, so its half-angle in \eqref{eq:lift} advances by only $\pi$,
landing on $-1$.  The longitude factor advances by $4\pi$, half-angle
$2\pi$, landing on $+1$ and contributing nothing.
\end{proof}

Proposition~\ref{prop:spinorial} is confirmed numerically three ways, all
returning $\sigma = -1.000000$: the adapted frame $[\bm{T},\bm{U},
\bm{T}\times\bm{U}]$ lifted by continuity, the framing field $A(\varphi)$
lifted by continuity, and direct evaluation of \eqref{eq:lift}.

\section{Stage 3: the value $2$ does no work}

Generalising Proposition~\ref{prop:spinorial} to the $(p,q)$ torus curve
with torus framing: the longitude half-angle advances by $\pi p$ and the
meridian half-angle by $\pi q$, so
\begin{equation}
\label{eq:parity}
\sigma = (-1)^{p+q}.
\end{equation}
Since $Sl = pq$ and $\gcd(p,q)=1$ forbids both being even, $p+q$ is odd
exactly when $pq$ is even, hence
\begin{equation}
\label{eq:parity2}
\sigma = (-1)^{Sl+1}.
\end{equation}

This was checked against an independent control family --- round circles
carrying $n$ framing twists, for which $Wr = 0$ and $|Sl| = n$ --- with
the measured result reproducing \eqref{eq:parity2} exactly:

\begin{center}
\begin{tabular}{@{}rrrrl@{}}
\toprule
$n$ & $Tw$ & $Wr$ & $Sl$ & class \\
\midrule
$-2$ & $+2$ & $0$ & $+2$ & spinorial \\
$-1$ & $+1$ & $0$ & $+1$ & trivial \\
$0$  & $0$  & $0$ & $0$  & spinorial \\
$1$  & $-1$ & $0$ & $-1$ & trivial \\
$2$  & $-2$ & $0$ & $-2$ & spinorial \\
$3$  & $-3$ & $0$ & $-3$ & trivial \\
$4$  & $-4$ & $0$ & $-4$ & spinorial \\
\bottomrule
\end{tabular}
\end{center}

\begin{corollary}[The identification $Sl=2 \leftrightarrow g=2$ loses its support]
\label{cor:kill}
The spin-relevant content of the self-linking number is one bit, its
parity.  A framed loop with $Sl = 0$ --- an ordinary round circle with an
untwisted framing --- is spinorial in precisely the same sense as the
M\"obius screw with $|Sl| = 2$.  The \emph{value} $2$ therefore carries no
spin information that $0$ does not also carry, and cannot be what
produces the double cover.
\end{corollary}

\begin{remark}[Status change]
The M\"obius-screw note (\S3.3) writes that $Sl = 2$ ``matches the $4\pi$
(two-turn) return of a spin-$1/2$ frame,'' and labels this a model
identification rather than a derivation.  Corollary~\ref{cor:kill}
removes the geometric motivation for the identification itself: the
matching quantity is a $\mathbb{Z}/2$ class, and $2$ and $0$ are the same
class.  Under the repository's tier discipline this is a
\textbf{T4 (explicitly falsified)} outcome for the reading of $Sl = 2$ as
the geometric origin of $g = 2$, and should be recorded in the
Elimination Ledger.  The rest of the M\"obius-screw note is untouched:
the geometry, the winding numbers, and the $|Sl| = 2$ computation are all
confirmed.
\end{remark}

\section{What survives, and it is worth keeping}

Something real remains, and it is more specific than what was claimed.
The double cover in this model is produced by the \textbf{meridian
half-angle} $\varphi/2$ appearing in \eqref{eq:curve} and
\eqref{eq:framing} --- that is, by the \emph{odd meridian winding}
$q = 1$, via \eqref{eq:parity}.  The screw's defining feature, that the
frame must traverse the curve twice to return, is a genuine topological
property of the shape, proved in Proposition~\ref{prop:spinorial}, and it
is exactly the $720^\circ$ property of a spin-$1/2$ frame.

So the model does land on the spinorial side of a genuine $\mathbb{Z}/2$
invariant.  That is a legitimate structural result and it is what the
tornado/Archimedean-screw picture was reaching for.  It is also
considerably weaker than a derivation of $g$: a $\mathbb{Z}/2$ bit cannot
produce a real number, and any model claiming a magnitude from this
structure needs a mechanism that \eqref{eq:parity} does not supply.

\begin{remark}[Falsifiable successor]
The next testable step is not a larger analogy but a smaller question:
does the framed-loop geometry generate a \emph{magnetic moment coupling},
by computing the current distribution on the ribbon and its magnetic
moment against the angular momentum?  That would produce a $g$ with
dimensions and a value, and could then fail.  The parity result above
neither supports nor blocks it.
\end{remark}

\section*{Reproduction}

\begin{verbatim}
python3 code/framed_unknot/framing_transformer.py
\end{verbatim}
Grid $N = 2000$ for the $(2,1)$ curve, $N = 1200$ for the controls,
$R = 1$, $a = 0.3$.  Results archived at
\texttt{docs/framed\_unknot\_results.json}.

\begin{thebibliography}{9}

\bibitem{calugareanu1961}
G.~C\u{a}lug\u{a}reanu,
\textit{Sur les classes d'isotopie des noeuds tridimensionnels et leurs invariants},
Czechoslovak Mathematical Journal \textbf{11} (1961), 588--625.

\bibitem{white1969}
J.~H.~White,
\textit{Self-linking and the Gauss integral in higher dimensions},
American Journal of Mathematics \textbf{91} (1969), 693--728.

\bibitem{fuller1971}
F.~B.~Fuller,
\textit{The writhing number of a space curve},
Proceedings of the National Academy of Sciences \textbf{68} (1971), 815--819.

\bibitem{mobiusscrew2026}
Flag Condensate Collaboration,
\textit{The M\"obius-Screw Electron: Framed Unknot Geometry for $g=2$ and a
Capacitance Model of $\alpha$},
preprint (July 2026),
\texttt{papers/notes/Mobius\_Screw\_Electron.tex}.

\end{thebibliography}

\end{document}
